\chapter{From fields to energies}
\label{ch:energies}

Everything up to this point produces \emph{fields}. This chapter closes the
loop: it assembles $\rho$, $\tau$ and $\vext$ into the Kohn--Sham total-energy
expression, states precisely which terms Poraqu\^e can and cannot evaluate, and
quantifies how accurate the fields must be before the resulting energy means
anything. The implementation is \pyclass{poraque.physics.energy}; the ASE
interface built on it is \pyclass{poraque.calculator.Poraque}.

\section{The expression}

\begin{equation}
  E = \underbrace{\int\tau\,d^3r}_{\Ts}
    + \underbrace{\int\rho\vext\,d^3r + E_{\alpha Z}}_{E_{\rm ext}}
    + \underbrace{\frac12\int\rho\vH\,d^3r}_{E_{\rm H}}
    + E_{\rm xc}[\rho]
    + E_{\rm Ewald}
  \label{eq:total-energy}
\end{equation}

Only the first term uses $\tau$, and it uses it in the simplest possible way:
$\Ts$ \emph{is} the integral of the kinetic energy density. That is the whole
reason \opt{chg2tau} is worth training --- it makes the one term of the
Kohn--Sham functional that has no closed form available from the density alone.

\subsection{Hartree}

Poisson in reciprocal space,
\begin{equation}
  \vH(\GG) = \frac{4\pi e^2\rho(\GG)}{G^2},
  \qquad \vH(\GG = 0) = 0 ,
\end{equation}
so $E_{\rm H} = \frac12\int\rho\vH\,d^3r$. Dropping $\GG = 0$ is the
neutralizing-background convention, chosen to match
\pyclass{ExternalPotential} so the two potentials are directly addable.

\begin{ptip}
The implementation is validated against an exact case rather than a recorded
output: for $\rho = \rho_0 + A\cos(\GG\cdot\rr)$ the energy is
$\pi e^2A^2\Omega/G^2$ in closed form, and the FFT result reproduces it to
$8\times10^{-16}$ relative. A uniform density gives exactly zero, which is the
correct answer in this convention rather than a numerical accident.
\end{ptip}

\subsection{Exchange and correlation}

Dirac exchange plus Perdew--Wang 1992 correlation, both local:
\begin{equation}
  \varepsilon_{\rm x} = -\frac34\left(\frac{3}{\pi}\right)^{1/3}\rho^{1/3},
  \qquad
  \varepsilon_{\rm c}(r_s) = -2A(1+\alpha_1 r_s)
    \ln\!\left[1 + \frac{1}{2A\,\Phi(r_s)}\right],
\end{equation}
with $\Phi(r_s) = \beta_1 r_s^{1/2} + \beta_2 r_s + \beta_3 r_s^{3/2}
+ \beta_4 r_s^{2}$ and $r_s = (3/4\pi\rho)^{1/3}$.

\begin{pwarn}
LDA is used regardless of the functional the reference calculation ran with,
and the mismatch is real: a PBE density evaluated with an LDA $E_{\rm xc}$ is
not a PBE energy. A gradient-corrected functional needs $\nabla\rho$ of a
\emph{predicted} field, whose noise the semilocal enhancement factor amplifies.
That deserves its own validation rather than a silent default.
\end{pwarn}

The density is clipped at zero first. Band-limiting a field with sharp core
peaks rings, so a resampled or predicted $\rho$ can undershoot slightly;
$\rho^{4/3}$ of a negative number is not real, and the physical content of
those points is ``no electrons here''.

\section{The \texorpdfstring{$\GG = 0$}{G = 0} bookkeeping}

This is the subtle part of Equation~\ref{eq:total-energy}, and getting it
wrong produces a plausible number that is wrong by eV per atom.

Each of the three electrostatic terms --- electron-ion, Hartree, ion-ion ---
diverges at $\GG = 0$, and in a neutral cell the three divergences cancel
exactly. The standard plane-wave treatment drops $\GG = 0$ from all three and
adds back the finite remainder of the electron-ion term. Writing VASP's local
potential as
\begin{equation}
  \vext(\GG) = \frac{1}{\Omega}\sum_s S_s(\GG)
    \left[v^{s}_{\rm short}(G) - \frac{4\pi Z^{\rm val}_s e^2}{G^2}\right],
\end{equation}
the divergent piece is the second bracketed term and the finite remainder at
$\GG = 0$ is $S_s(0) = N_s$ times $v^{s}_{\rm short}(0)$. Hence
\begin{equation}
  E_{\alpha Z} = \frac{N_{\rm elec}}{\Omega}\sum_s N_s\,\mathrm{PSCORE}_s,
  \qquad \mathrm{PSCORE}_s = \lim_{q\to0}v^{s}_{\rm short}(q),
\end{equation}
which is exactly VASP's \opt{alpha Z}. It is read straight from the
\file{POTCAR} tables --- the same tables Chapter~\ref{ch:vasp} uses for the
potential itself --- via
\opt{PotcarSingle.pscore}.

\begin{pnote}
When the tables are unavailable the term is reported as \opt{None}, and
\opt{EnergyComponents.missing} names it. Defaulting to zero would look like a
working calculation while silently dropping a term of order eV per atom, which
is precisely the size of the effects one would want to resolve.
\end{pnote}

\section{Ewald}

The ion-ion energy of point charges $q_i$ in a neutralizing background:
\begin{align}
  E_{\rm Ewald} = {}& \frac{e^2}{2}\sideset{}{'}\sum_{i,j,\Rv}
      \frac{q_iq_j\,\mathrm{erfc}(\eta|\rr_{ij}+\Rv|)}{|\rr_{ij}+\Rv|}
    + \frac{2\pi e^2}{\Omega}\sum_{\GG\neq0}
      \frac{e^{-G^2/4\eta^2}}{G^2}\bigl|S(\GG)\bigr|^2 \nonumber\\
   & - \frac{\eta e^2}{\sqrt\pi}\sum_i q_i^2
     - \frac{\pi e^2}{2\eta^2\Omega}\Bigl(\sum_i q_i\Bigr)^2 .
\end{align}

The last term is the interaction with the compensating background. It is not
optional: without it the result depends on the arbitrary splitting parameter
$\eta$ instead of being invariant under it --- which is also the cheapest
available test that the implementation is correct.

\begin{ptip}
Verified against tabulated Madelung constants: rocksalt $1.747564595$ to
$8\times10^{-13}$, caesium chloride $1.762674773$ to $7\times10^{-11}$,
zincblende $1.6380550$ to $5\times10^{-8}$; and against the jellium limit
$E = -\xi q^2e^2/2L$, $\xi = 2.8372974795$, for a single ion in a cube. Each
depends on every part of the sum being individually right.
\end{ptip}

\section{What the number means}
\label{sec:energy-meaning}

\subsection{It is not a DFT total energy}

\file{CHGCAR} holds the valence \emph{pseudo} density and \file{TAUCAR} the
valence pseudo kinetic energy density. The PAW one-centre terms --- the
difference between the pseudo and all-electron descriptions inside the
augmentation spheres --- never appear in any file Poraqu\^e reads, so they are
absent from Equation~\ref{eq:total-energy} entirely. For the 27-atom Au
supercell the result is $\approx -31215$ eV against a VASP \opt{TOTEN} of order
$-100$ eV. No adjustment of the terms above closes that gap; the missing
physics is simply not in the inputs.

\subsection{The cancellation problem}

The individual terms are large and of alternating sign:

\begin{center}
\begin{tabular}{@{}lr@{}}
\toprule
Term & eV \\
\midrule
$\Ts$          & $+6\,207$ \\
$E_{\rm ext}$  & $-12\,635$ \\
$E_{\alpha Z}$ & $+1\,753$ \\
$E_{\rm H}$    & $+3\,230$ \\
$E_{\rm xc}$   & $-3\,821$ \\
$E_{\rm Ewald}$ & $-25\,949$ \\
\midrule
Total & $-31\,215$ \\
\bottomrule
\end{tabular}
\end{center}

Across the seven reference supercells the \emph{physically relevant} variation
of that total is $7.9$ eV --- a relative $2.5\times10^{-4}$ of its own
magnitude. That sets a brutal accuracy requirement on the fields, because every
electrostatic term is linear or quadratic in $\rho$.

Measured with the current universal models, evaluated on their own training
data (so this is a lower bound on the true error):

\begin{center}
\begin{tabular}{@{}lr@{}}
\toprule
Quantity & Value \\
\midrule
True spread of $E$ across the seven structures & $7.9$ eV \\
MAE of the predicted energy differences & $22.3$ eV ($0.83$ eV/atom) \\
Correlation of predicted vs.\ true differences & $r = 0.61$ \\
\bottomrule
\end{tabular}
\end{center}

\begin{pwarn}
The error is roughly \textbf{three times the signal}. Energy differences from
this pipeline are not yet usable. The cause is cancellation, not a defect in
the energy module: a field-level relative $L^{2}$ of $3\times10^{-2}$ cannot
survive a $2.5\times10^{-4}$ signal-to-total ratio. Reaching $0.01$ eV/atom
requires densities accurate to roughly $10^{-5}$ relative --- three orders of
magnitude beyond the current models.
\end{pwarn}

Two candidate explanations were tested and rejected:

\begin{description}
  \item[Grid resolution.] On reference fields the total energy shifts by
    $0.08$ eV out of $31\,215$ ($0.003$ eV/atom) between \opt{resolution: 32}
    and the native $128^3$ mesh, and both $\Ts$ and the electron count are
    preserved \emph{exactly} --- spectral resampling is an exact band-limited
    projection. The grid is not the limiting factor.
  \item[Electron-count drift.] Predicted densities miss the exact valence count
    by up to $2.8\,\%$. Rescaling $\rho$ to the known total makes the result
    \emph{worse} ($3.20$ eV/atom, from $0.83$), because $E_{\rm ext}$ is linear
    in $\rho$ and a $0.3\,\%$ rescale shifts a $-12\,690$ eV term by $38$ eV.
    The drift is a symptom, not the cause.
\end{description}

\subsection{Consequences for the programme}

This is the sharpest quantitative target the project has produced so far, and
it reframes what ``accurate enough'' means. A relative $L^{2}$ of $0.03$ on the
density is an excellent \emph{field} prediction and a useless \emph{energy}
one. Two directions follow, and they are not equivalent:

\begin{enumerate}
  \item \textbf{Make the fields three orders of magnitude better.} Direct, and
    almost certainly unreachable by scaling the current architecture.
  \item \textbf{Make the energy insensitive to the field error.} Predict the
    energy \emph{difference} from a reference state rather than the absolute
    total, or exploit the variational property of the Kohn--Sham functional:
    at the true density the energy is stationary, so a first-order density
    error produces only a second-order energy error. That property is
    available only if the energy is evaluated with a \emph{self-consistent}
    functional, which is what the physics-informed programme of
    Chapter~\ref{ch:pifno} is ultimately for.
\end{enumerate}

The second route is why $\delta \Ts/\delta\rho$ (Chapter~\ref{ch:model2}) is
worth having: it is the ingredient that turns a learned $\tau$ into a
variational object rather than a fitted one.

\section{The ASE interface}

\pyclass{poraque.calculator.Poraque} attaches the whole chain to an
\opt{ase.Atoms} object:

\begin{lstlisting}
from ase.build import bulk
from poraque.calculator import Poraque

atoms = bulk("Au", "fcc", a=4.08, cubic=True)
atoms.calc = Poraque("models/ext2chg.pt", "models/chg2tau.pt",
                     potcar="POTCAR")
energy = atoms.get_potential_energy()
\end{lstlisting}

Each call runs \opt{Atoms} $\to$ \opt{Structure} $\to$ \opt{FieldGrid} $\to$
$\vext$ $\to$ $\rho$ $\to$ $\tau$ $\to$ $E$, and leaves the three fields on the
calculator so a prediction can be written out rather than only reduced to a
number.

It resembles MACE or NequIP at the interface but not underneath.
\opt{get\_forces()} and \opt{get\_stress()} raise \opt{NotImplementedError}, so
geometry optimisation and molecular dynamics are out of reach. Two independent
pieces are missing: $\partial\vext/\partial\Rv_a$ (analytic, a
Hellmann--Feynman term) and back-propagation of $\partial E/\partial\rho$
through both operators. A finite-difference stand-in is no shortcut --- on a
fixed grid it would be dominated by the grid's own discontinuity as atoms cross
voxel boundaries.

\begin{pwarn}
Without a \file{POTCAR} the calculator falls back to the Gaussian pseudo-ion
model of Chapter~\ref{ch:vasp}, whose relative $L^{2}$ against the tabulated
potential is $0.13$ --- four orders of magnitude worse than the
$2\times10^{-5}$ the operators were trained on. The input is then outside the
training distribution and the prediction is not meaningful. The calculator
warns once and proceeds, because the fallback remains useful for smoke tests.
\end{pwarn}
